
\chapter{Persistent and Distributed Transactional Memory}
Two extensions of the TM contract change the runtime far more than the
language: making commits survive crashes (persistent TM) and making the
same API span multiple machines (distributed TM). Both keep the ordinary
bank invariant, then add a stronger one --- durability, or agreement across
nodes --- and both force the commit path to become explicit about ordering.
The two worked examples below implement and verify each extension.

\section{Worked Example: Durable Bank Commit}
\index{persistent TM}\index{durability}\index{redo log}\index{crash recovery}
\label{sec:durable-bank-example}

Persistent TM asks for the ordinary TM invariant plus a crash invariant. A
committed transfer must not disappear after power loss, and recovery must not
apply it twice. The small example below uses one durable intent record. The
record stores the post-state, so redo is idempotent.

\begin{lstlisting}[language=C, caption={A minimal durable two-account transfer with idempotent redo.}]
#include <array>
#include <cstdint>
#include <cstddef>
#include <stdexcept>
#include <fcntl.h>
#include <unistd.h>

static const uint64_t MAGIC = 0x544d5f42414e4b31ULL; /* "TM_BANK1" */

struct Journal {
    uint64_t magic;
    uint32_t from, to;
    int64_t old_from, old_to;
    int64_t new_from, new_to;
    uint64_t sum;
};

struct Image {
    Journal journal;
    std::array<int64_t, 2> balance;
    int64_t floor;
};

static uint64_t checksum(const Journal& j) {
    return j.magic ^ j.from ^ (uint64_t(j.to) << 7) ^
           uint64_t(j.old_from) ^ (uint64_t(j.old_to) << 11) ^
           uint64_t(j.new_from) ^ (uint64_t(j.new_to) << 17);
}

static bool valid(const Journal& j) {
    return j.magic == MAGIC && j.sum == checksum(j);
}

static void write_all(int fd, const void* p, size_t n, off_t off) {
    const char* b = static_cast<const char*>(p);
    while (n != 0) {
        ssize_t k = ::pwrite(fd, b, n, off);
        if (k <= 0) throw std::runtime_error("pwrite failed");
        b += k; n -= size_t(k); off += k;
    }
}

static void read_all(int fd, void* p, size_t n, off_t off) {
    char* b = static_cast<char*>(p);
    while (n != 0) {
        ssize_t k = ::pread(fd, b, n, off);
        if (k <= 0) throw std::runtime_error("pread failed");
        b += k; n -= size_t(k); off += k;
    }
}

class DurableBank {
    int fd_;
    Image img_{};

    void sync() {
        if (::fsync(fd_) != 0) throw std::runtime_error("fsync failed");
    }

    void clear_journal() {
        img_.journal = Journal{};
        write_all(fd_, &img_.journal, sizeof(Journal), 0);
        sync();
    }

    void flush_image() {
        write_all(fd_, &img_, sizeof(Image), 0);
        sync();
    }

    void recover() {
        if (!valid(img_.journal)) return;
        const Journal& j = img_.journal;
        if (j.from >= 2 || j.to >= 2) throw std::runtime_error("bad journal");
        img_.balance[j.from] = j.new_from;  /* idempotent redo */
        img_.balance[j.to] = j.new_to;
        write_all(fd_, &img_.balance, sizeof(img_.balance),
                  offsetof(Image, balance));
        sync();
        clear_journal();
    }

public:
    explicit DurableBank(const char* path) {
        fd_ = ::open(path, O_RDWR | O_CREAT, 0644);
        if (fd_ < 0) throw std::runtime_error("open failed");

        off_t n = ::lseek(fd_, 0, SEEK_END);
        if (n < off_t(sizeof(Image))) {
            img_.balance = {1000, 1000};
            img_.floor = 0;
            flush_image();
        } else {
            read_all(fd_, &img_, sizeof(Image), 0);
            recover();
        }
    }

    ~DurableBank() {
        if (fd_ >= 0) ::close(fd_);
    }

    bool transfer(uint32_t from, uint32_t to, int64_t amount) {
        if (from >= 2 || to >= 2 || from == to || amount < 0) return false;
        if (img_.balance[from] - amount < img_.floor) return false;

        Journal j{};
        j.magic = MAGIC;
        j.from = from; j.to = to;
        j.old_from = img_.balance[from];
        j.old_to = img_.balance[to];
        j.new_from = j.old_from - amount;
        j.new_to = j.old_to + amount;
        j.sum = checksum(j);

        img_.journal = j;
        write_all(fd_, &img_.journal, sizeof(Journal), 0);
        sync();                         /* durable commit intent */

        img_.balance[from] = j.new_from;
        img_.balance[to] = j.new_to;
        write_all(fd_, &img_.balance, sizeof(img_.balance),
                  offsetof(Image, balance));
        sync();                         /* durable data */

        clear_journal();                /* recovery no longer needed */
        return true;
    }

    int64_t total() const {
        return img_.balance[0] + img_.balance[1];
    }
};

int main() {
    DurableBank b("bank.img");
    b.transfer(0, 1, 5);
    return b.total() == 2000 ? 0 : 2;
}
\end{lstlisting}

\begin{itemize}
    \item The durable intent record is not a full database log. It is the minimum needed to make one committed transfer recoverable.
    \item The post-state makes recovery idempotent: replaying the same journal record writes the same balances.
    \item The account floor is checked before the durable commit point. Recovery must not re-run policy code; it must complete the already chosen outcome.
    \item The money-conservation invariant is checked after restart by the same observation used in volatile tests: the total balance remains unchanged.
\end{itemize}

\section{Worked Example: Model Checking Distributed Commit}
\index{distributed TM}\index{two-phase commit}\index{TLA+}\index{money conservation}
\label{sec:distributed-commit-example}

Distributed TM turns validation into a protocol. A coordinator may decide only
after participants vote, and the data update must remain atomic from the
client's point of view. This TLA+ model is intentionally small: two accounts,
one transfer, one prepare phase, and one commit decision.

\begin{lstlisting}[caption={A TLA+ model of one distributed bank transfer with prepare/commit.}]
--------------------------- MODULE MiniDistributedBank ---------------------------
EXTENDS Naturals, TLC

Accounts == 1..2
From == 1
To == 2
Amount == 3
Floor == 0
InitBal == [a \in Accounts |-> 10]

VARIABLES bal, phase, vote, decision

vars == <<bal, phase, vote, decision>>

Init ==
    /\ bal = InitBal
    /\ phase = "idle"
    /\ vote = [r \in Accounts |-> "none"]
    /\ decision = "none"

Begin ==
    /\ phase = "idle"
    /\ phase' = "prepare"
    /\ UNCHANGED <<bal, vote, decision>>

Prepare(r) ==
    /\ phase = "prepare"
    /\ vote[r] = "none"
    /\ vote' =
        [vote EXCEPT ![r] =
            IF r = From
            THEN IF bal[From] - Amount >= Floor THEN "yes" ELSE "no"
            ELSE "yes"]
    /\ UNCHANGED <<bal, phase, decision>>

AllVoted ==
    \A r \in Accounts: vote[r] \in {"yes", "no"}

DecideCommit ==
    /\ phase = "prepare"
    /\ \A r \in Accounts: vote[r] = "yes"
    /\ phase' = "commit"
    /\ decision' = "commit"
    /\ UNCHANGED <<bal, vote>>

DecideAbort ==
    /\ phase = "prepare"
    /\ AllVoted
    /\ \E r \in Accounts: vote[r] = "no"
    /\ phase' = "done"
    /\ decision' = "abort"
    /\ UNCHANGED <<bal, vote>>

ApplyCommit ==
    /\ phase = "commit"
    /\ bal' = [bal EXCEPT ![From] = @ - Amount,
                          ![To] = @ + Amount]
    /\ phase' = "done"
    /\ UNCHANGED <<vote, decision>>

Next ==
    \/ Begin
    \/ \E r \in Accounts: Prepare(r)
    \/ DecideCommit
    \/ DecideAbort
    \/ ApplyCommit

Spec == Init /\ [][Next]_vars

TypeOK ==
    /\ phase \in {"idle", "prepare", "commit", "done"}
    /\ decision \in {"none", "commit", "abort"}
    /\ \A a \in Accounts: bal[a] \in Nat
    /\ \A r \in Accounts: vote[r] \in {"none", "yes", "no"}

MoneyConserved ==
    bal[1] + bal[2] = InitBal[1] + InitBal[2]

FloorsHold ==
    \A a \in Accounts: bal[a] >= Floor

=============================================================================
\end{lstlisting}

\begin{itemize}
    \item The model separates decision from application. This exposes the window in which a real coordinator must recover enough state to finish.
    \item The invariant is the same as the sequential bank example: balances stay above the floor and the total amount of money is conserved.
    \item The prepare votes are local, but the commit is global. This is the essential cost of exporting a TM abstraction over a distributed store.
    \item The model is small enough for exhaustive checking, but already contains the main distributed failure boundary: a decision without a completed data update.
\end{itemize}

\section{Summary}
\begin{itemize}
    \item Persistent TM adds a crash invariant to the ordinary money-conservation invariant; the durable worked example shows the minimum journal needed.
    \item Distributed TM turns validation into a protocol; the model-checking example pins the failure boundary at ``a decision without a completed data update.''
    \item Both extensions keep the same TM API --- only the commit path and its ordering obligations change.
\end{itemize}

\section{Exercises}
\begin{enumerate}
    \item \emph{[implementation]} Extend the durable bank example in \S\ref{sec:durable-bank-example} to four accounts and two concurrent worker threads. Keep one global durable commit lock, then record the throughput change relative to an in-memory SGL version. Do not claim durability until the journal and data writes have both been forced to storage.
    \item \emph{[verification]} Modify the TLA+ model in \S\ref{sec:distributed-commit-example} so that the coordinator may crash after deciding commit but before applying the update. Add a recovery action and check that \texttt{MoneyConserved} and \texttt{FloorsHold} still hold.
\end{enumerate}

%  PART VIII: HORIZONS
